% ---------------------------------------------------------------------------
% The CDM-frame division of the integral solution, term by term.
%
% SELF-CONTAINED BY DESIGN. This file inputs nothing and defines nothing; it
% can be moved into Paper 1 by moving one file. Decision S42's pattern.
%
% This is the bundle's contribution to Paper 1's work-list item 3 --- *"The
% CDM-frame division has not been computed term by term"* --- and it is
% offered for checking by the Paper 1 stream, not as a settled result. Every
% step whose content is a sign or an ell-weight is flagged as such.
% ---------------------------------------------------------------------------
\section{The CDM-frame division of the integral solution, term by term}
\label{sec:frame-division}

Paper~1's Sec.~III\,C records that \emph{``the CDM-frame division has not been
computed term by term''}. This appendix computes it for the integral solution
(102), finds that it does not close, and localises the failure to a single
coefficient at a single power of $k$. \textbf{It is offered for checking rather
than as a settled result}: the three transformation laws it rests on are
transcriptions, and under the rule that a claim whose content is a sign or an
$\ell$-weight is settled in the index algebra or not settled, the index algebra
has not been done here.

\subsection{What has to hold}

Let $n^a$ be the shear-free (Newtonian) frame and $u^a=n^a+v^a$ the
zero-acceleration (CDM) frame, with $v_an^a=0$ and $v^a$ first order. A
first-order change of frame is a boost, and a boost changes the observed
brightness by $\tau\to\tau-v_ae^a$: \textbf{a pure dipole}. In PSTF multipoles,
\begin{equation}
  \tau^{(u)}_{A_1}=\tau^{(n)}_{A_1}-v_{A_1},
  \qquad
  \tau^{(u)}_{A_\ell}=\tau^{(n)}_{A_\ell}\quad\text{for all }\ell\ge2 .
  \label{eq:boost-is-a-dipole}
\end{equation}
\eqref{eq:boost-is-a-dipole} is not an approximation to be checked against a
tolerance. \textbf{It is what makes $C_\ell$ an observable at all}, and the
almost-EGS argument of \emph{Annals~II}'s Appendix~E rests on it. So the
division below has only one admissible outcome: the shift must cancel
identically for $\ell\ge2$.

\subsection{The three shifts}

The coefficients are fixed by (90),
\begin{equation}
  -C_0\simeq\tfrac52\,a\mathcal B_2,\qquad
  -C_1\simeq\tfrac1k\big(a\mathcal B_1+\kappa'v_B\big),\qquad
  -C_2\simeq\tfrac1{k^2}\,a\mathcal B_2 ,
  \label{eq:C-from-B}
\end{equation}
with $a\mathcal B_1=k\,\delta T+aA$ and $v_B=k\tilde\tau_1$. The frame
carries four entries, each from the antecedent:

\begin{center}
\begin{tabular}{lll}
\hline
 & Newtonian $n^a$ & CDM $u^a$ \\
\hline
$aA$ & $k\Phi_A$ & $0$ \quad(the frame condition $A_a=0$) \\
$\delta T$ & $\delta\tilde T$ & $\delta\tilde T-\mathcal Hv/k$ \quad(n-t3) \\
$\tilde\tau_1$ & $\tilde\tau_1$ & $\tilde\tau_1-v/k$ \quad(n-t4) \\
$a\mathcal B_2$ & $0$ & $-kv$ \quad(n-t1, $a\sigma=-kv$ exactly) \\
\hline
\end{tabular}
\end{center}

\noindent
Write $\Delta X\equiv X^{(u)}-X^{(n)}$. Then
$\Delta(a\mathcal B_1)=-\mathcal Hv-k\Phi_A$, and the Euler equation
$v'+\mathcal Hv=-k\Phi_A$ --- which \emph{is} the frame condition $A_a=0$ ---
collapses it to $v'$. The other two follow directly:
\begin{equation}
  \boxed{\;
  \Delta C_0=+\tfrac52\,k\,v,\qquad
  \Delta C_1=-\tfrac1k\big(v'-\kappa'v\big),\qquad
  \Delta C_2=+\tfrac{v}{k}\;}
  \label{eq:three-shifts}
\end{equation}

\textbf{This shows that} the three coefficients pick up the frame velocity at
three different powers of $k$: $C_0$ at $k^{+1}$, $C_1$ at $k^{-1}$, $C_2$ at
$k^{-1}$. That spread is not a choice --- it is \eqref{eq:C-from-B}, where
$C_0$ and $C_2$ are \emph{the same} $a\mathcal B_2$ divided by $k^0$ and $k^2$.

\subsection{The shift of (102), term by term}

(102) reads
\begin{equation}
  \frac{\beta_\ell\tau_\ell(\eta_0)}{2\ell+1}\simeq
  \big[C_2'(\eta_*)-C_1(\eta_*)\big]\mathcal D\,j_\ell(x_*)
  +kC_2(\eta_*)\mathcal D\,j_\ell'(x_*)
  +\int_{\eta_*}^{\eta_0}\!\!d\eta\;\mathcal V e^{-(k/k_D)^2}
   \big[C_0-C_1'+C_2''\big]j_\ell(x).
  \label{eq:102-here}
\end{equation}
Applying \eqref{eq:three-shifts} to each bracket:
\begin{align}
  \Delta\big[C_2'-C_1\big]_{\eta_*} &= \tfrac1k\big(2v'-\kappa'v\big)\Big|_{\eta_*},
  \label{eq:shift-scalar}\\[2pt]
  \Delta\big[kC_2\big]_{\eta_*} &= v(\eta_*),
  \label{eq:shift-shear}\\[2pt]
  \Delta\big[C_0-C_1'+C_2''\big] &= \tfrac52kv
    +\tfrac1k\big(2v''-(\kappa'v)'\big).
  \label{eq:shift-integrand}
\end{align}

\textbf{The first term of \eqref{eq:shift-integrand} has no partner.} Count
powers of $k$ against the rate at which $v$ evolves. $v$ obeys
$v'=-\mathcal Hv-k\Phi_A$, so on sub-horizon scales $v$ changes on the Hubble
rate, not on $k$: $v'\sim\mathcal Hv$ and $v''\sim\mathcal H^2v$. Then
\begin{equation}
  \frac{\tfrac52kv}{\tfrac1k\,2v''}\;\sim\;\frac54\left(\frac{k}{\mathcal H}\right)^{2},
  \label{eq:no-partner}
\end{equation}
and the two boundary shifts \eqref{eq:shift-scalar}--\eqref{eq:shift-shear} are
$O(\mathcal Hv/k)$ and $O(v)$, smaller still. \textbf{Nothing in (102) is the
same size as $\tfrac52kv$}, so the shift cannot cancel, and
\eqref{eq:boost-is-a-dipole} fails for every $\ell\ge2$ by a factor that grows
as $(k/\mathcal H)^2$.

\subsection{Measured}

Computed in both frames from one environment, one acoustic solution and one
$k$ grid --- the only difference being which frame (110) specialises to ---
against the Newtonian total's own scale over $\ell\in\{2,5,10,50,200\}$:

\begin{center}
\begin{tabular}{lrrr}
\hline
background & $k$ & as assembled & with $C_0$ dropped \\
\hline
CDM (Einstein--de Sitter) & 20  & $53.5$   & $0.032$ \\
                          & 300 & $1719$   & $0.035$ \\
$\Lambda$CDM              & 20  & $62.8$   & $0.105$ \\
                          & 300 & $6868$   & $0.154$ \\
\texttt{PLANCK\_MODEL}    & 20  & $72.2$   & $0.111$ \\
                          & 300 & $3687$   & $0.058$ \\
\hline
\end{tabular}
\end{center}

\noindent
Three to four orders of magnitude, to a residual of $1$--$15\%$, and the $k$
dependence is the $(k/\mathcal H)^2$ of \eqref{eq:no-partner}. The rank
decomposition says the same from the other side: $C_0$ carries $98.6\%$ of
$C_\ell$ in the CDM frame and is identically zero in the Newtonian one.

\subsection{Why \emph{Annals~II} never meets this}

(102) is derived for a general $u^a$ --- the text says so explicitly after
(90), \emph{``Note that these are CGI with respect to $u^a$''} --- and then
evaluated once, in the Newtonian frame, where \emph{``the vanishing shear
condition $\tilde\sigma_{ab}\approx0$ implies that $\tilde C_0(\eta)\approx0$
and $\tilde C_2(\eta)\approx0$''}. \textbf{Both members of the pair vanish
together}, so the imbalance between them cannot be seen. Nothing the antecedent
publishes is affected. What is missing is a stated condition on a formula it
labels frame-general, and it is missing because the paper had no reason to
carry (102) anywhere the condition fails.

\subsection{Three repairs tested and rejected}

Recorded because a rejected hypothesis that leaves no trace is formed again.

\paragraph{The dropped lines of (93).} (102) is (93) with its last two
integrals removed, \emph{``we can effectively drop the last two lines''}; both
carry $C_1$ and $C_2$, so they are the obvious counterweight. Restoring
$\int(\kappa'e^{-\kappa})[C_1-2C_2]j_\ell$ and
$\int((\kappa')^2-\kappa'')e^{-\kappa}C_2j_\ell$ moves the disagreement from
$72.2$ to $71.6$ at $k=20$ and from $3686.9$ to $3686.8$ at $k=300$. They are
not the counterweight.

\paragraph{Undoing the integration by parts.} (93) comes from the kernel
$A_\ell=C_0\tau^{(0)}_\ell+C_1\partial_{\eta'}\tau^{(0)}_\ell
+C_2\partial^2_{\eta'}\tau^{(0)}_\ell$, where the two derivatives act on the
\emph{kernel} and bring down $k$, not on the coefficient where they bring down
$\mathcal H$. That suggests the by-parts step is what strands $C_0$.
\textbf{The numerical test of this was invalid} --- the implementation failed
its own control, reproducing the by-parts form in the Newtonian frame only to
within factors of $10^{-3}$ to $2.5$ --- because (102) does not differentiate
the weight $\mathcal V$ that (93) carries inside the integral, so the two forms
are not related by an exact by-parts once a weight is present. \textbf{The idea
is not refuted; it is untested}, and it is the first thing Paper~1 should try.

\paragraph{Reading the $\ell=0$ line of (89) as mis-sourced.} $a\mathcal
B_0\delta_{\ell0}$ appears in the full hierarchy, is commented out of the
equation the integral solution solves, and is restored in the open-case
extension --- which looks like a dropped source. It is not. The ansatz must
produce no monopole, $\tau_0=0$, and differentiating $\int A_0\,d\eta'=0$ gives
exactly (89)'s $\ell=0$ line with \textbf{zero} on the right. $C_0$ is
correctly derived.

\subsection{What Paper 1 should check, and in what order}

\begin{enumerate}
\item \textbf{The four entries of the frame table}, by index algebra. $a\sigma=-kv$
  (n-t1), $\delta T\to\delta T-\mathcal Hv/k$ (n-t3) and
  $\tilde\tau_1\to\tilde\tau_1-v/k$ (n-t4) are transcriptions in this bundle.
  They carry the signs that \eqref{eq:three-shifts} is built from, and a sign
  is settled in the index algebra or it is not settled.
\item \textbf{\eqref{eq:boost-is-a-dipole} at second order.} The division here
  is first order. The $1$--$15\%$ residual left after $C_0$ is removed is the
  honest size of what is not yet accounted for, and whether it is
  $O(\varepsilon^2)$ or an artefact of the grids is open.
\item \textbf{The by-parts step}, as above --- the one untested repair.
\item \textbf{Whether $C_0$ belongs in the integrand}, which is this bundle's
  question~Q14 and which the measurement answers empirically and the algebra
  does not.
\end{enumerate}

\textbf{What this appendix establishes} is narrower than any of those and does
not depend on them: \emph{whatever the resolution, the CDM-frame division of
(102) does not close, the failure is in $C_0$ alone, and its size is
$(k/\mathcal H)^2$.} That is one coefficient and one power of $k$, where
Paper~1's Sec.~III\,C had an open gap.
